%-------------------------------------------------------------------------------
%	SECTION TITLE
%-------------------------------------------------------------------------------
\cvsection{Experiencia Laboral}


%-------------------------------------------------------------------------------
%	CONTENT
%-------------------------------------------------------------------------------
\begin{cventries}
%---------------------------------------------------------
  \cventry
    {Ingeniero de Software Senior - MSAI (Search, Assistant \& Intelligence)} % Job title
    {Microsoft} % Organization
    {Santiago, Chile} % Location
    {\daterange{1}{2024}{9}{2025}} % Date(s)
    {
        \begin{cvitems}
            % --- CDAC ---
            \item {
                Diseñé la arquitectura de CDAC (Certification Detection and Control), un sistema distribuido de alerta temprana creado para modernizar el pipeline de aseguramiento de calidad de los productos de búsqueda de M365 (Word, PowerPoint, Outlook, SharePoint, etc).
                  \begin{itemize}
                    \item {Revolucioné el proceso de pruebas de regresión reemplazando ciclos de recolección de telemetría A/B de semanas por técnicas avanzadas de análisis estadístico.}
                    \item {Reduje el tiempo de detección de regresiones de métricas a tan solo 6 horas, evitando reescrituras de código y resguardando la calidad de búsqueda mientras se cumplían las agresivas demandas de despliegue de la era Copilot.}
                  \end{itemize}
            }
            % --- LSS ---
            \item {
              Desarrollé funcionalidades core de LSS (Local Shard Search), un microservicio de alta concurrencia que sirve como el puente crítico entre las APIs de búsqueda y los shards de almacenamiento globales, manejando más de 200M de requests diarios.
              \begin{itemize}
                \item {Lideré el desarrollo de capacidades de búsqueda exclusivas para Microsoft Copilot, habilitando acceso a nuevos índices de datos y patrones de consulta avanzados que enriquecieron significativamente las respuestas generadas por IA para análisis de documentos complejos.}
                \item {Desarrollé agentes autónomos internos basados en LLMs para analizar telemetría y diagnosticar incidentes, reduciendo la carga operativa manual del on-call en aproximadamente 50\% y agilizando el mantenimiento del sistema.}
                \item {Optimicé componentes críticos de baja latencia migrando código legacy de C\# a Rust, mejorando significativamente el throughput y la eficiencia de recursos.}
                \item {Mantuve alta disponibilidad en infraestructura de búsqueda crítica para el negocio, realizando rotaciones de on-call en un entorno de producción a gran escala.}
              \end{itemize}
            }
            % --- STACK ---
            \item {Tecnologías utilizadas: C\#, .NET, Rust, Python, KQL, Azure, Herramientas internas de LLMs.}
        \end{cvitems}
    }
%---------------------------------------------------------
  \cventry
    {Director de Ingeniería LATAM} % Job title
    {Nezasa} % Organization
    {Santiago, Chile} % Location
    {\daterange{9}{2022}{9}{2023}} % Date(s)
    {
      \begin{cvitems} % Description(s) of tasks/responsibilities
        \item {Lideré la fusión de las áreas tecnológicas entre TripYeah y Nezasa, asegurando una integración eficiente.}
        \item {Coordiné y planifiqué actividades para equipos en toda LATAM, optimizando recursos y plazos.}
        \item {Dirigí la integración de TripYeah (ahora TripOptimizer) en la suite de productos de Nezasa, mejorando significativamente el proceso de creación de itinerarios.}
        \item {Implementé un chatbot innovador usando la API de OpenAI, mejorando la interfaz de creación de itinerarios de viaje.}
        \item {Gestioné la resolución de incidentes y el desarrollo de funcionalidades avanzadas.}
        \item {Tecnologías utilizadas: TypeScript, ReactJS, NextJS, Golang, Python, Kubernetes, Skaffold, Kafka, gRPC, Protobuf, AWS, APIs de OpenAI.}
      \end{cvitems}
    }

%---------------------------------------------------------
  \cventry
    {Co-fundador \& CTO} % Job title
    {TripYeah [Adquirida por Nezasa]} % Organization
    {Santiago, Chile} % Location
    {\daterange{4}{2021}{9}{2022}} % Date(s)
    {
      \begin{cvitems} % Description(s) of tasks/responsibilities
        \item {Coordiné desarrollos entre equipos distribuidos en Latinoamérica.}
        \item {Diseñé una arquitectura de microservicios event-driven para resolver problemas complejos de itinerarios multi-destino, integrando precios de vuelos en tiempo real de proveedores como Amadeus o Sabre.}
        \item {Mejoré significativamente la eficiencia de las recomendaciones (itinerarios multi-destino con precios en tiempo real), reduciendo los tiempos de 90 a 10 segundos usando patrones como CQRS y arquitectura event-driven.}
        \item {Desarrollé soluciones Full Stack y lideré la resolución de incidentes y desarrollos complejos.}
        \item {Mentoreé ingenieros trainee y contribuí al levantamiento de capital (+2M USD) y la venta de la compañía.}
        \item {Tecnologías utilizadas: TypeScript, ReactJS, NextJS, Golang, Python, Kubernetes, Skaffold, Kafka, gRPC, Protobuf, Redis, AWS.}
      \end{cvitems}
    }
\end{cventries}
\newpage
\begin{cventries}

%---------------------------------------------------------
  \cventry
    {Data Scientist y Tech Leader} % Job title
    {Globant} % Organization
    {Santiago, Chile} % Location
    {\daterange{11}{2019}{4}{2021}} % Date(s)
    {
      \begin{cvitems} % Description(s) of tasks/responsibilities
        \item {Lideré el equipo de Augmented Design dentro de la iniciativa Augmented Globant (ahora Globant X), enfocado en el desarrollo de una plataforma innovadora de pruebas de usabilidad UX. Esta plataforma se distinguió por su capacidad de capturar datos faciales, tono de voz y seguimiento ocular, habilitando estadísticas avanzadas de interacción de usuarios.}
        \item {Responsable de planificar, crear e implementar modelos avanzados de machine learning y deep learning, especializados en reconocimiento de expresiones faciales y emociones en tonos de voz, así como en la generación de heatmaps mediante eye tracking.}
        \item {Colaboré estrechamente con el Product Manager en la planificación del desarrollo de software, asegurando la alineación con los objetivos del proyecto.}
        \item {Desempeñé un rol clave como Full Stack Developer, contribuyendo al desarrollo integral de la plataforma.}
        \item {Participé en presentaciones y pre-ventas en Brasil y Chile, representando al estudio de data science de Globant.}
        \item {Obtuve una categorización directa de Semi-Senior a Senior (saltando Semi-Senior Advance) durante el primer año, gracias a buenas evaluaciones y gestión de proyectos.}
        \item {Tecnologías utilizadas: TypeScript, ReactJS, Golang, Python, TensorFlow, TF.JS, Docker, Docker-compose, AWS.}
      \end{cvitems}
    }

%---------------------------------------------------------
  \cventry
    {Co-fundador \& Full Stack Developer} % Job title
    {Ticketplus.cl} % Organization
    {Santiago, Chile} % Location
    {\daterange{3}{2015}{11}{2017}} % Date(s)
    {
      \begin{cvitems} % Description(s) of tasks/responsibilities
        \item {Desarrollé y operé una plataforma innovadora de venta de tickets para eventos, optimizando la experiencia de usuario y la eficiencia del proceso de venta.}
        \item {Interacción directa con clientes para la mejora y personalización del producto.}
        \item {Lideré la planificación y ejecución de nuevas features y funcionalidades, mejorando significativamente la oferta de servicios de la plataforma.}
        \item {Implementé un sistema eficiente de validación en puerta de eventos, asegurando una gestión de acceso ágil y segura.}
        \item {Diseñé e introduje un sistema de cola online para atención de compras masivas, mejorando la capacidad de manejo de alto tráfico y ventas.}
        \item {Tecnologías utilizadas: JavaScript, AngularJS, Ruby on Rails, Redis, AWS.}
      \end{cvitems}
    }

\end{cventries}
